% Monthly Notices of the Royal Astronomical Society.
\documentclass[fleqn,usenatbib]{mnras}

\usepackage[T1]{fontenc}
\usepackage{graphicx}
\usepackage{amsmath}

\title[An Example Title]{An Example Title}
\author[A. U. Thor et al.]{
A. U. Thor,$^{1}$\thanks{E-mail: author@example.org}
and Second Author$^{2}$
\\
$^{1}$Institute of Astrophysics, Your University, City, Country\\
$^{2}$Another Institute, City, Country
}
\date{Accepted XXX. Received YYY; in original form ZZZ}
\pubyear{2026}

\begin{document}
\label{firstpage}
\pagerange{\pageref{firstpage}--\pageref{lastpage}}
\maketitle

\begin{abstract}
State the question, the data and method, and the main result. MNRAS limits
the abstract to 250 words and asks for a single paragraph.
\end{abstract}

\begin{keywords}
galaxies: distances and redshifts -- stars: variables: general
\end{keywords}

\section{Introduction}

Say what is known, what is missing, and what this paper adds. Cite with
\citet{einstein1905} in the text or in parentheses \citep{knuth1984}.

\section{Methods}

Describe the observations and the analysis in enough detail to repeat them.

\begin{equation}
  \mu = 5 \log_{10}\left(\frac{d}{10\,\mathrm{pc}}\right)
  \label{eq:distance}
\end{equation}

\begin{table}
  \centering
  \caption{Example measurements.}
  \label{tab:results}
  \begin{tabular}{lcc}
    \hline
    Object & $\mu$ (mag) & $\sigma_\mu$ (mag) \\
    \hline
    A & 18.48 & 0.03 \\
    B & 18.52 & 0.04 \\
    \hline
  \end{tabular}
\end{table}

\section{Conclusions}

Give the result and what follows from it.

\section*{Acknowledgements}

Thank the people and grants that made the work possible.

\section*{Data Availability}

Say where the data can be found. MNRAS requires this section.

\bibliographystyle{mnras}
\bibliography{references}

\bsp
\label{lastpage}
\end{document}
